\graphicspath{{conclusion/fig/}}

\mychapter{Conclusion}{Conclusion}{}\label{chap:conclusion}

\mysection{Summary}{Summary}\label{sec:sum}

The use of dynamic time warping (DTW) as a Bryde's whale call detection system was investigated in this project. Three feature 
extraction methods were tested: short-time Fourier transform (STFT), mel-frequency cepstral coefficients (MFCC) 
and linear predictive coding coefficients (LPCC).

\mysection{Objectives}{Objectives}\label{sec:obj}
The first objective of this project was to implement a \acs{DTW}-based detection system using three feature extraction methods: 
the short-time Fourier transform (STFT), mel-frequency cepstral coefficients (MFCCs) and linear predictive coding coefficients (LPCC). 
This was designed and implemented as explained in Chapter~\ref{chap:sys-design}.

The second objective was to tune the feature frame length, call type thresholds and vote fraction for each method based on \gls{f1_score}. 
The results for this can be seen in Tables~\ref{tab:stft_calib_results},~\ref{tab:stft_vote_results},~\ref{tab:mfcc_calib_results},~\ref{tab:mfcc_vote_results},~\ref{tab:lpc_calib_results} 
and~\ref{tab:lpc_vote_results}. 

The third objective was to evaluate the detection performance of each configuration using \gls{precision}, \gls{recall}, \gls{f1_score} 
and average precision (\acs{AP}). This can be seen in Tables~\ref{tab:stft_test_results},~\ref{tab:mfcc_test_results} and~\ref{tab:lpc_test_results}.

\mysection{Limitations}{Limitations}\label{sec:lim}

Limitations in this project included the scarcity of well-labelled data and the human error involved in the 
labelling, due to a range of human labellers being used for different sections. This meant that it was necessary 
to select call type and background calibration templates, which caused more human error within the template 
selection stage. The necessity of calibration template selection led to an imbalance between calls and background 
noise in hyperparameter tuning and evaluation. The labelling based on previous projects also led to many false positives (FPs) 
actually being calls that were never labelled.

\mysection{Conclusion}{Conclusion}\label{sec:con}

The STFT model performed the best out of the three considered models. The maximum f1-score across all 
STFT experiments was 0.792 and the maximum AP was 0.860. The LPCC model performed second best with best scores of 
0.650 and 0.789 for f1-score and AP respectively. The MFCC model performed the worst with a highest f1-score and AP value of 
0.647 and 0.684 respectively. 

For the MFCC experiments the addition of derivatives did not improve the model. For all three models an 
increase in template diversification (5 vs 10 templates per call type) had little effect on hyperparameter selection. 
It was found that the single-tone thresholds chosen for all experiments were generally better than the other 
thresholds, with a possible link to the simplicity of their frequency structure in comparison to the other 
two call types. Multi-tone call thresholds always had the lowest f1-scores, probably due to human error caused 
by different definitions decided by each labeller.

Precision was consistently lower than recall, showing that all models were more lenient in the evaluation 
stage than intended during the hyperparameter selection stage. Inspection of FP calls for the STFT model 
found that many of the FP calls were actually unlabelled calls. This shows that the limitations may have caused 
a skewed precision result in the evaluation analysis. The presence of FP calls closer to already labelled calls and in the second 
half of the 72-hour file pointed to possible connections between SNR and call density and FP detections. 

Overall the STFT model was the most suitable method by a sizeable margin in both f1-score and AP values for detection of 
Bryde's whale calls in the False Bay region.

\mysection{Future Work}{Future Work}\label{sec:fut_work}

More consistent labelling and more in-depth labelling of all audio files in future would be beneficial. A model 
could be tuned on full audio files instead of selected templates, and many of the limitations would be removed. 
The consideration of different methods to remove background noise would also be beneficial.

Once the limitations of the current detection model have been addressed, the next step would be to 
extend the system to classify the call type of each detection in continuous audio files.